%! AP Statistics — Unit 1, Lesson 1.10 — Student Packet Cover Sheet
\documentclass[10pt]{article}
\usepackage{apstats-article}
\usepackage{apstats-boxes}

%=============================================================================
\begin{document}

% ── Full-bleed navy banner ────────────────────────────────────────────────────
\begin{tikzpicture}[remember picture, overlay]
  \fill[navy] (current page.north west)
    rectangle ([yshift=-0.9in]current page.north east);
\end{tikzpicture}
\vspace{-0.8in}

\begin{center}
  {\color{white}\textbf{\LARGE AP Statistics}}\\[4pt]
  {\color{skymid}\large\textbf{Unit 1: Exploring One-Variable Data}}\\[6pt]
  {\color{white}\textbf{\large Lesson 1.10 \quad The Normal Distribution and $z$-Scores}}
\end{center}

\vspace{0.22in}

\namedateperiod

\vspace{0.18in}

% ── LEARNING TARGETS ─────────────────────────────────────────────────────────
\begin{learningtargetbox}
\small
\begin{enumerate}[leftmargin=1.5em, itemsep=5pt]
  \item \ldots{} describe the shape of the Normal distribution and identify
        $\mu$ and $\sigma$ as its two parameters, using the notation $N(\mu, \sigma)$.
  \item \ldots{} apply the \textbf{68--95--99.7 (Empirical) Rule} to find approximate
        proportions of values within 1, 2, or 3 standard deviations of the mean.
  \item \ldots{} calculate a $z$-score using $z = \dfrac{x - \mu}{\sigma}$, interpret
        it in context, and use the Normal table or technology to find probabilities
        and percentiles.
  \item \ldots{} assess whether a distribution is approximately Normal using a
        histogram or Normal probability plot.
\end{enumerate}
\end{learningtargetbox}

\vspace{0.14in}

% ── TABLE OF CONTENTS ────────────────────────────────────────────────────────
\begin{tocbox}
\small
\renewcommand{\arraystretch}{1.5}
\begin{tabularx}{\linewidth}{@{} c l X r @{}}
\rowcolor{sky}
\textbf{\#} & \textbf{Component} & \textbf{Description} & \textbf{Score} \\
\midrule
1 & Warm-Up        & Read a Normal curve; identify mean, SD, and the 68\% region            & \blank{1.2cm} \\
2 & Guided Notes   & Normal distribution properties; Empirical Rule; $z$-scores; Normal table & \blank{1.2cm} \\
3 & Group Activity & Apply Empirical Rule; compute $z$-scores; use Normal table with SAT data & \blank{1.2cm} \\
4 & Exit Ticket    & Find $z$-score, probability, and percentile for female heights           & \blank{1.2cm} \\
5 & Homework       & Empirical Rule; $z$-scores and table; AP-style free response              & \blank{1.2cm} \\
\midrule
  & \multicolumn{2}{r}{\textbf{Total}} & \blank{1.2cm} \\
\end{tabularx}
\end{tocbox}

\vspace{0.14in}

% ── KEEP IN MIND ─────────────────────────────────────────────────────────────
\begin{remindbox}
\small
Lessons 1.6--1.9 described and compared distributions using graphs and summary
statistics. Lesson 1.10 introduces the \textit{Normal model} --- a mathematical
distribution that lets us calculate precise probabilities when the shape is
symmetric and bell-shaped.

\vspace{6pt}
\begin{tabularx}{\linewidth}{@{} >{\centering\arraybackslash\columncolor{goldbg}}p{0.06\linewidth}
                                    X @{}}
\textbf{\large!} &
\textbf{The Normal table gives area to the LEFT.}\enspace
To find a ``greater than'' probability, subtract from 1.
To find an ``between'' probability, subtract the smaller area from the larger. \\[6pt]
\textbf{\large!} &
\textbf{Always convert $x$ to a $z$-score first.}\enspace
The formula $z = \dfrac{x - \mu}{\sigma}$ standardizes any Normal variable.
Never look up a raw $x$-value in the Normal table. \\[6pt]
\textbf{\large!} &
\textbf{The Empirical Rule is an approximation.}\enspace
68--95--99.7\% are rounded values for exactly 1, 2, and 3 standard deviations.
Use the Normal table (or technology) when the problem asks for a precise probability. \\[6pt]
\textbf{\large!} &
\textbf{Real data is never perfectly Normal.}\enspace
Always check the histogram or Normal probability plot before applying the Normal model.
The model is an approximation --- use it when it is reasonable, not automatically.
\end{tabularx}
\end{remindbox}

\end{document}
